\documentclass[11pt]{article}

\usepackage[a4paper,margin=1in]{geometry}
\usepackage[T1]{fontenc}
\usepackage[utf8]{inputenc}
\usepackage{lmodern}
\usepackage{amsmath,amssymb,amsthm,mathtools}
\usepackage{bm}
\usepackage[dvipsnames]{xcolor}
\usepackage{enumitem}
\usepackage{hyperref}
\usepackage[nameinlink,noabbrev]{cleveref}

\hypersetup{
  colorlinks=true,
  linkcolor=blue!50!black,
  citecolor=blue!50!black,
  urlcolor=blue!50!black,
}

\newtheorem{theorem}{Theorem}[section]
\newtheorem{proposition}[theorem]{Proposition}
\newtheorem{corollary}[theorem]{Corollary}
\newtheorem{lemma}[theorem]{Lemma}
\theoremstyle{definition}
\newtheorem{definition}[theorem]{Definition}
\newtheorem{remark}[theorem]{Remark}
\newtheorem{assumption}[theorem]{Assumption}

\newcommand{\R}{\mathbb{R}}
\newcommand{\E}{\mathbb{E}}
\newcommand{\PP}{\mathcal{P}}
\newcommand{\W}{W_1}
\newcommand{\dd}{\,\mathrm{d}}

\title{Mean Field Theory for BatchEnsembles}
\author{}
\date{\today}

\begin{document}
\maketitle

\begin{abstract}
BatchEnsemble introduces diversity through multiplicative fast weights, but not all such modulations play the same role when width grows. This note studies a fixed-$K$, width-to-infinity mean-field limit for a vector-valued bias-free two-layer BatchEnsemble-style model with neuron-wise fast weights and optional boundary fast weights. The main conclusion is structural. Neuron-wise modulations are averaged out by width, whereas boundary fast weights are not. As a result, in the regime without boundary fast weights, the infinite-width dynamics reduce to a single deterministic transport PDE and permutation-symmetric initialization forces exact head collapse. In contrast, when at least one boundary fast weight is retained, the correct large-width object is a deterministic PDE-ODE system for the neuron law and the boundary state, and order-one head dependence can survive through the finite-dimensional initialization of that state. We prove the supporting well-posedness, stability, and finite-width convergence results in an abstract bounded-Lipschitz setting and then specialize them to these two practitioner-relevant regimes.
\end{abstract}

\section{Introduction}

BatchEnsemble is a parameter-sharing ensemble architecture: instead of training
$K$ fully independent networks, one keeps a shared backbone and gives each head
a small set of multiplicative fast weights. The practical promise is clear. One
would like to obtain ensemble-like diversity without paying the full cost of $K$
separate models. The theoretical question is less clear: \emph{which
head-specific parameters still matter once width becomes very large?}

This note studies that question in a fixed-$K$, width-$M\to\infty$ mean-field
regime. The key point is that BatchEnsemble mixes two different kinds of
head-specific structure:
\begin{enumerate}[leftmargin=2em]
  \item \emph{Neuron-wise fast weights}: variables such as
  $s_{\alpha,m}^{(2)}$ and $r_{\alpha,m}^{(1)}$ that are attached to individual
  neurons and therefore
  self-average as $M\to\infty$.
  \item \emph{Boundary fast weights}: variables such as $s_\alpha^{(1)}$ and
  $r_\alpha^{(2)}$ that sit at the input or output boundary, stay
  finite-dimensional when $M\to\infty$, and therefore do not self-average.
\end{enumerate}
The distinction is operational. Width averaging acts directly on neuron-wise
modulations, but it does not remove asymmetry carried by boundary fast weights. In the concrete
two-layer model introduced in \Cref{sec:prelim}, this leads to a sharp split
between two regimes. If all boundary fast weights are fixed to the all-ones
vectors, then every head-specific quantity lives inside the neuron state and the
width limit reduces to a single transport PDE. If at least one boundary fast
factor is retained, then the correct large-width limit is a coupled PDE-ODE
system and order-one head dependence can survive through the finite-dimensional
boundary-state initialization.

The main message is practitioner-facing. If head diversity is implemented only
through neuron-wise modulations, then width averaging removes it at leading
order: under standard permutation-symmetric initialization, all heads collapse
to the same infinite-width predictor. If one wants order-one head asymmetry to
survive width averaging, then a boundary fast weight must be kept in the state.
In the full model, the directly input- and output-connected factors
$s_\alpha^{(1)}$ and $r_\alpha^{(2)}$ play that role. Once such a factor is
retained, the correct limit is no longer a single PDE but a coupled PDE-ODE
system. This
point is especially notable when compared with NTK-based analyses of embedded
ensembles, where one can obtain an \emph{independent regime} under centered
last-layer post-activation modulation \cite{velikanov2022}. Our claim is
different: in the present feature-learning mean-field regime, neuron-wise
modulations alone do not generically produce leading-order functional
independence.

The rest of the note is organized as follows. \Cref{sec:prelim} introduces the
concrete two-layer model and the neuron-wise versus boundary-fast-weight split.
\Cref{sec:related} positions the paper relative to BatchEnsemble, NTK-based
embedded-ensemble theory, feature-learning mean-field theory, and the tabular
model TabM. \Cref{sec:results} then states the two main implications for
BatchEnsemble design, compares them with the NTK-based regime picture, and
records the technical fixed-$K$ width limit that supports both
specializations. The abstract mean-field setup, standing assumptions, and all
proofs are deferred to the appendix.

\section{Related Work}
\label{sec:related}

\paragraph{BatchEnsemble and parameter-efficient ensembling.}
BatchEnsemble was introduced by Wen et al.\ \cite{wen2020} as a
parameter-sharing alternative to deep ensembles, in which each ensemble member
is represented by rank-one multiplicative fast weights applied to shared weight
matrices. Our concrete two-layer model is a mean-field idealization of this
construction. The present note is not about improving the architecture itself,
but about understanding which parts of this parameter-efficient ensemble remain
visible in the wide feature-learning limit.

\paragraph{NTK-based operating regimes for embedded ensembles.}
The closest theoretical comparison is the NTK-based analysis of Embedded
Ensembles by Velikanov et al.\ \cite{velikanov2022}. In their lazy-training
setting, infinite-width dynamics are described by a deterministic ensemble NTK,
and centered last-layer post-activation modulation can produce an
\emph{independent regime} in which off-diagonal GP covariances and NTK entries
vanish. Our main message differs precisely at this point: in the
feature-learning mean-field regime studied here, neuron-wise modulation alone
does not generically create leading-order functional independence. Instead, the
law-of-large-numbers limit collapses unless some boundary fast weight survives
width averaging.

\paragraph{Feature-learning mean-field theory.}
On the mean-field side, two-layer networks admit a deterministic
distributional description in terms of transport PDEs or Wasserstein gradient
flows; see, for example, Mei et al.\ \cite{mei2019} and the lazy-training
perspective of Chizat et al.\ \cite{chizat2019}. Our contribution is not a new
generic mean-field theorem for ordinary networks, but rather an application of
the feature-learning mean-field viewpoint to parameter-efficient ensembles with
shared weights. The outcome is a structural distinction between neuron-wise
fast weights, which self-average with width, and boundary fast weights, which
remain as finite-dimensional state variables.

\paragraph{TabM and the tabular case for efficient ensembling.}
This distinction matters in practice because BatchEnsemble-style ensembling has
already become influential in tabular deep learning. TabM, introduced by
Gorishniy et al.\ \cite{gorishniy2025}, is an MLP-based tabular model built
around parameter-efficient ensembling and reports a strong
performance-efficiency trade-off on public benchmarks. Our results do not
analyze TabM directly, but they help articulate why boundary-connected
modulation mechanisms can be qualitatively different from purely neuron-wise
ones in the feature-learning regime. In that sense, the present theory is
meant as a step toward understanding why parameter-efficient ensembling can be
so effective in domains such as tabular learning.

\section{Preliminary}
\label{sec:prelim}

\subsection{Concrete two-layer BatchEnsemble model}

Fix an input dimension $d_x$, an output dimension $d_y$, a hidden width $M$, and an ensemble size
$K\in\{1,2,\dots\}$. For each neuron $m\in\{1,\dots,M\}$, let
\[
  a_m\in\R^{d_y},\qquad
  w_m\in\R^{d_x},\qquad
  s_{\alpha,m}^{(2)}\in\R\quad (\alpha=1,\dots,K),\qquad
  r_{\alpha,m}^{(1)}\in\R\quad (\alpha=1,\dots,K),
\]
and let
\[
  s^{(1)}=(s_1^{(1)},\dots,s_K^{(1)})\in(\R^{d_x})^K,
  \qquad
  r^{(2)}=(r_1^{(2)},\dots,r_K^{(2)})\in(\R^{d_y})^K.
\]
The two layers take the standard BatchEnsemble form
\[
  W_\alpha = W\circ \bigl(s_\alpha^{(1)}(r_\alpha^{(1)})^\top\bigr),
  \qquad
  A_\alpha = A\circ \bigl(s_\alpha^{(2)}(r_\alpha^{(2)})^\top\bigr),
\]
and the corresponding bias-free two-layer predictor is
\begin{equation}
  f_{\alpha,\mathrm{full}}^M(x)
  :=
  \frac1M A_\alpha^\top \sigma(W_\alpha^\top x).
  \label{eq:model-full-matrix}
\end{equation}
Equivalently, expanding row by row,
\begin{equation}
  f_{\alpha,\mathrm{full}}^M(x)
  :=
  \frac1M\sum_{m=1}^M
  (a_m\odot r_\alpha^{(2)})\,s_{\alpha,m}^{(2)}\,
  \sigma\!\bigl(
    r_{\alpha,m}^{(1)}\,w_m^\top(s_\alpha^{(1)}\odot x)
  \bigr).
  \label{eq:model-full}
\end{equation}
The scalar output case is recovered by taking $d_y=1$.

The width-$M$ empirical measure is taken only over the neuron-indexed
coordinates. That is,
\begin{equation}
  \rho^M
  :=
  \frac1M\sum_{m=1}^M \delta_{\theta_m},
  \qquad
  \theta_m=
  (a_m,w_m,s_{1,m}^{(2)},\dots,s_{K,m}^{(2)},r_{1,m}^{(1)},\dots,r_{K,m}^{(1)}).
  \label{eq:empirical}
\end{equation}
The boundary fast weights are kept outside this empirical measure:
\begin{equation}
  \bar h^M := (s^{(1)},r^{(2)}) \in (\R^{d_x})^K\times(\R^{d_y})^K.
  \label{eq:boundary-concrete}
\end{equation}
Thus the full finite-width state is the pair $(\rho^M,\bar h^M)$, and the
head-specific variables split naturally into
\begin{enumerate}[leftmargin=2em]
  \item neuron-level coordinates
  $(s_{1,m}^{(2)},\dots,s_{K,m}^{(2)},r_{1,m}^{(1)},\dots,r_{K,m}^{(1)})$
  inside $\theta_m$,
  \item the boundary fast weights collected in $\bar h^M=(s^{(1)},r^{(2)})$.
\end{enumerate}
This concrete split already identifies the objects that appear in the
large-width limit: a neuron law $\rho_t$ and, when boundary fast weights are
present, a finite-dimensional boundary state $h_t$. The rigorous abstract
mean-field formulation and the standing regularity assumptions are recorded in
Appendix~\ref{sec:appendix-setup}; the main text uses only the resulting
interpretation.

\section{Theoretical Results}
\label{sec:results}

\paragraph{Mean-field setup in one sentence.}
The technical results below are stated in terms of the abstract mean-field
objects introduced in Appendix~\ref{sec:appendix-setup}. In that notation, the
large-width limit is described by a neuron law $\rho_t$ and, when boundary fast
weights are present, a finite-dimensional boundary state $h_t$; the full
feature-learning limit is a coupled PDE-ODE system \eqref{eq:pde}--\eqref{eq:ode},
while the no-boundary regime removes the ODE and leaves a single transport PDE.

\paragraph{Guiding question.}
The central design question is which head-specific parameters survive width
averaging. The results below show a sharp split. Neuron-wise modulations
$(s_{\alpha,m}^{(2)},r_{\alpha,m}^{(1)})$ are averaged into the neuron law, while
boundary fast weights remain finite-dimensional state variables. In the full
two-layer model these are $s_\alpha^{(1)}$ and $r_\alpha^{(2)}$. This
structural split is what drives the different practical interpretations of the
two regimes.

\subsection{Without boundary fast weights, wide BatchEnsemble collapses}
\label{sec:regime-a}

In the concrete full model \eqref{eq:model-full-matrix}--\eqref{eq:model-full},
the regime without boundary fast weights is obtained by fixing
\[
  s_\alpha^{(1)}\equiv \mathbf 1,
  \qquad
  r_\alpha^{(2)}\equiv \mathbf 1
  \qquad
  \text{for all }\alpha.
\]
Then the predictor reduces to
\begin{equation}
  f_{\alpha,\mathrm{free}}^M(x)
  :=
  \frac1M\sum_{m=1}^M
  a_m\,s_{\alpha,m}^{(2)}\,
  \sigma\!\bigl(
    r_{\alpha,m}^{(1)}\,w_m^\top x
  \bigr).
  \label{eq:model-free}
\end{equation}
At the abstract level, this is exactly the specialization in which the feature
maps are independent of the boundary state. Formally, this means
\[
  \Psi_\alpha(x;\theta,h)\equiv \Psi_\alpha^{\mathrm{free}}(x;\theta),
\]
which includes \eqref{eq:model-free}. Then \eqref{eq:drift-G} vanishes and the
ODE for $h_t$ disappears.

\begin{corollary}[No boundary fast weights reduce the limit to a single PDE]
\label{cor:boundaryfree}
Assume \cref{ass:main} and suppose the feature maps are independent of $h$. Then
the fixed-$K$ mean-field system reduces to the transport PDE
\begin{equation}
  \partial_t \rho_t
  +
  \nabla_\theta\cdot\bigl(\rho_t B(\theta;\rho_t)\bigr)
  =
  0,
  \label{eq:pde-free}
\end{equation}
where
  \[
  B(\theta;\rho)
  :=
  -\E_{(x,y)\sim\mathsf P}
  \Bigl[
    \frac1K\sum_{\alpha=1}^K
    \nabla_\theta\Psi_\alpha^{\mathrm{free}}(x;\theta)^\top
    \nabla_1\ell(f_\alpha(x;\rho),y)
  \Bigr].
\]
All conclusions of \cref{thm:deterministic,cor:random} remain valid with the
boundary state deleted.
\end{corollary}

In the regime without boundary fast weights, standard i.i.d.\ initialization from a
deterministic law produces a deterministic infinite-width limit because there is
no residual boundary-state random environment.

\begin{corollary}[Without boundary fast weights there is no order-one width-limit randomness]
\label{cor:boundaryfree-det}
In the setting of \cref{cor:boundaryfree}, assume that the empirical initial laws
$\rho_0^M$ converge to a deterministic law $\rho_0$ in probability, almost
surely, or in $L^1$. Then the infinite-width predictor vector
\[
  \bigl(f_1(\cdot;\rho_t),\dots,f_K(\cdot;\rho_t)\bigr)
\]
is deterministic.
\end{corollary}

The decisive simplification is that, once the boundary state is removed, all
head dependence lives inside the same neuron law.

\begin{proposition}[Permutation-symmetric initialization forces head collapse]
\label{prop:collapse}
Assume the no-boundary setting of \cref{cor:boundaryfree}. If the initial law
$\rho_0$ is invariant under every $\Pi_\pi^\Theta$, i.e.
\[
  (\Pi_\pi^\Theta)_\#\rho_0=\rho_0
  \qquad
  \forall \pi\in S_K,
\]
then the unique characteristic solution of \eqref{eq:pde-free} satisfies
\[
  (\Pi_\pi^\Theta)_\#\rho_t=\rho_t
  \qquad
  \forall t\ge 0,\ \forall \pi\in S_K.
\]
Consequently,
\begin{equation}
  f_1(x;\rho_t)=\cdots=f_K(x;\rho_t)
  \qquad
  \forall x\in\mathcal X,\ \forall t\ge 0.
  \label{eq:collapse}
\end{equation}
\end{proposition}

\begin{corollary}[Standard symmetric initialization collapses at leading order]
\label{cor:boundaryfree-collapse}
In the setting of \cref{cor:boundaryfree-det}, if the deterministic limit law
$\rho_0$ is permutation-symmetric in the head coordinates, then the infinite-width
limit satisfies \eqref{eq:collapse}.
\end{corollary}

\begin{remark}[Contrast with NTK-based operating regimes]
\label{rem:ntk-contrast}
This is the sharpest point of contrast with NTK-based analyses of embedded
ensembles. In the lazy-training picture of \cite{velikanov2022}, centered
neuron-wise post-activation modulation can produce an \emph{independent regime}
in which off-diagonal GP covariances and NTK entries vanish. By contrast, the
present feature-learning mean-field limit shows that, once all boundary fast
weights are removed, permutation-symmetric initialization yields exact head
collapse at leading order. Thus NTK-style independence does not arise here as a
law-of-large-numbers effect of neuron-wise modulation alone; if it appears at
all, it belongs to a lazy or fluctuation-level description rather than to the
leading-order feature-learning limit.
\end{remark}

Taken together, \cref{cor:boundaryfree,cor:boundaryfree-det,prop:collapse,cor:boundaryfree-collapse}
say that neuron-wise fast weights alone do not create leading-order functional
diversity in the fixed-$K$, wide-limit regime. For a practitioner, this means
that removing all boundary fast weights pushes BatchEnsemble toward a
single-model mean-field limit, even though finite-width heads may still differ through
subleading fluctuations.

\subsection{With boundary fast weights, asymmetry survives width averaging}
\label{sec:regime-b}

In the concrete full model \eqref{eq:model-full-matrix}--\eqref{eq:model-full},
this means retaining at least one of the directly input- and output-connected
factors $s_\alpha^{(1)}$ or $r_\alpha^{(2)}$. Once a boundary fast weight is
kept, the width limit remains deterministic
conditional on $(\rho_0,h_0)$, but need not be deterministic unconditionally.

\begin{corollary}[Surviving order-one randomness comes from boundary fast weights]
\label{cor:randomness}
Assume \cref{ass:main}. If $(\rho_0,h_0)$ is deterministic, then the infinite-width
predictor vector is deterministic. If $\rho_0$ is deterministic but $h_0$ is
random, then any order-one randomness that survives at infinite width is induced
entirely by the law of $h_0$.
\end{corollary}

In this regime, permutation symmetry of the neuron law alone no longer forces
collapse, because the predictors also depend on boundary coordinates. To
obtain collapse one would need a stronger symmetric manifold that couples both
the neuron law and the boundary state. Without such a condition, the PDE-ODE
limit can remain head-dependent even at infinite width.

\begin{remark}[Comparison with the no-boundary regime]
The contrast with \cref{prop:collapse} is structural. Without boundary fast
modulation, all head dependence lives inside the neuron law, so permutation
symmetry of that law is enough to identify all head predictors. Once boundary
fast weights are present, the predictor map depends on an additional
finite-dimensional variable that width does not average out. This is why the
boundary regime naturally leads to a conditional deterministic theory rather
than to automatic leading-order collapse.
\end{remark}

\subsection{Implications for BatchEnsemble design}

The single formulation above makes the two regimes easy to compare.

\paragraph{State variable.}
Without boundary fast weights, the correct large-width state is just the neuron
law $\rho_t$. With boundary fast weights, it is the pair $(\rho_t,h_t)$.

\paragraph{Limit equation.}
Without boundary fast weights, one gets a single transport PDE. With
boundary fast weights, one gets a coupled PDE-ODE system.

\paragraph{Randomness at infinite width.}
Without boundary fast weights, standard deterministic-law initialization
produces a deterministic limit. With boundary fast weights, order-one
randomness can survive through the initial law of the finite-dimensional
boundary state.

\paragraph{Symmetry and collapse.}
Without boundary fast weights, permutation-symmetric initialization forces
exact head collapse in the infinite-width limit. With boundary fast weights,
symmetry of the neuron law alone is no longer enough; one must also control the
boundary state.

\paragraph{Interpretation.}
The no-boundary regime is the cleanest width mean-field limit for fixed $K$:
width averages everything that is head-specific. The boundary regime is the
next level of complexity: width still averages neuron-level fluctuations, but it
does not remove asymmetry carried by boundary fast weights.

\paragraph{Operational takeaway.}
If the goal is to preserve order-one head diversity at very large width, then
a boundary fast weight is structurally different from neuron-wise fast
weights and should not be viewed as an inessential implementation detail. In the
full model, that role is played by the directly input- and output-connected
factors $s^{(1)}$ and $r^{(2)}$. Conversely, if one removes all boundary fast
modulation, then the wide-limit theory predicts that standard symmetric
initialization produces a collapsed ensemble at leading order.

\subsection{Comparison with NTK-based operating regimes}
\label{sec:ntk-comparison}

The most relevant comparison is with the NTK-based analysis of Embedded
Ensembles by Velikanov et al.\ \cite{velikanov2022}. In that lazy-training
picture, the wide-network limit is described by the initialization GP and by a
deterministic ensemble NTK. Their main structural conclusion is that centered
post-activation modulation in the last hidden layer, namely
$U_1^L=\E[u_{\alpha,j}^L]=0$, yields an \emph{independent regime}: the
off-diagonal GP covariances and off-diagonal NTK entries vanish, so the
embedded ensemble behaves like an ensemble of independent reference networks.
When $U_1^L\neq 0$, they obtain a \emph{collective regime}, where these
off-diagonal quantities are nonzero.

Our fixed-$K$ mean-field theory makes a different statement because it studies a
different leading-order object. We are not linearizing around initialization.
Instead, we follow the feature-learning transport of the neuron law. In this
regime, neuron-wise modulations such as
$\bigl(s_{\alpha,m}^{(2)},r_{\alpha,m}^{(1)}\bigr)$ are absorbed into the common
law $\rho_t$ as $M\to\infty$. Consequently, in the no-boundary regime,
standard permutation-symmetric initialization leads to
\[
  f_1(\cdot;\rho_t)=\cdots=f_K(\cdot;\rho_t),
\]
so there is no analogue of a leading-order \emph{independent regime} generated
solely by centered neuron-wise modulations. In the full model, the only
head-specific mechanisms that survive width averaging at order one are the
boundary fast weights $s^{(1)}$ and $r^{(2)}$.

This difference is not a contradiction with the NTK theory. It indicates that
the two asymptotic descriptions live at different levels. The NTK independent
regime is a lazy, kernel-level notion of decoupling. Our result says that in a
feature-learning mean-field scaling, the law-of-large-numbers limit collapses
unless a non-self-averaging boundary variable is present. From this
perspective, NTK-style independence should be interpreted here as a property of
the lazy/fluctuation picture rather than as the generic leading-order behavior
of wide feature-learning BatchEnsembles.

For practitioners, this comparison suggests a concrete design lesson. If the
goal is to retain diversity in a regime where feature learning remains active,
then centered neuron-wise modulations are not enough by themselves. One needs a
mechanism that survives width averaging, namely some boundary fast weight, or
else the leading-order wide-limit theory predicts collapse.

\subsection{Technical backbone: general fixed-\texorpdfstring{$K$}{K} width limit}
\label{sec:technical-backbone}

The previous subsections isolate the practitioner-facing implications. Their
common technical input is the following fixed-$K$ width-limit theorem, which
turns the finite-width BatchEnsemble dynamics into a deterministic mean-field
system on finite horizons.

\begin{theorem}[Deterministic well-posedness, stability, and particle convergence]
\label{thm:deterministic}
Assume \cref{ass:main}, fix $T>0$, and let $(\rho_0,h_0)\in\PP_1(\Theta)\times\mathcal H$.

\begin{enumerate}[label=(\roman*),leftmargin=2em]
  \item There exists a unique characteristic solution
  \[
    (\rho_t,h_t)_{t\in[0,T]}
    \in
    C\bigl([0,T],\PP_1(\Theta)\times\mathcal H\bigr)
  \]
  of \cref{eq:pde,eq:ode}.
  \item There exists a constant $C_T<\infty$ such that for any two initial data
  $(\rho_0,h_0)$ and $(\mu_0,\bar h_0)$ with corresponding solutions
  $(\rho_t,h_t)$ and $(\mu_t,\bar h_t)$,
  \begin{equation}
    \sup_{t\in[0,T]}
    \bigl(
      \W(\rho_t,\mu_t)+\|h_t-\bar h_t\|
    \bigr)
    \le
    C_T
    \bigl(
      \W(\rho_0,\mu_0)+\|h_0-\bar h_0\|
    \bigr).
    \label{eq:det-stability}
  \end{equation}
  \item Let deterministic particles and boundary states satisfy
  \begin{align*}
    \dot\theta_{m,t}^M &= B(\theta_{m,t}^M;\rho_t^M,h_t^M),
    \qquad
    \rho_t^M=\frac1M\sum_{m=1}^M\delta_{\theta_{m,t}^M},\\
    \dot h_t^M &= G(\rho_t^M,h_t^M).
  \end{align*}
  If
  \[
    \W(\rho_0^M,\rho_0)+\|h_0^M-h_0\|\xrightarrow[M\to\infty]{}0,
  \]
  then
  \begin{equation}
    \sup_{t\in[0,T]}
    \bigl(
      \W(\rho_t^M,\rho_t)+\|h_t^M-h_t\|
    \bigr)
    \le
    C_T\bigl(\W(\rho_0^M,\rho_0)+\|h_0^M-h_0\|\bigr)
    \xrightarrow[M\to\infty]{}0.
    \label{eq:particle-conv}
  \end{equation}
  In particular,
  \[
    \sup_{t\in[0,T]}
    \|f_\alpha(x;\rho_t^M,h_t^M)-f_\alpha(x;\rho_t,h_t)\|
    \to 0
  \]
  for every fixed $\alpha$ and $x$.
\end{enumerate}
\end{theorem}

\begin{corollary}[Random initialization]
\label{cor:random}
Assume \cref{ass:main}, fix $T>0$, and let $(\rho_0,h_0)$ be a possibly random
initial environment. Suppose the random initial data $(\rho_0^M,h_0^M)$ satisfy
\[
  \W(\rho_0^M,\rho_0)+\|h_0^M-h_0\|\xrightarrow[M\to\infty]{}0
\]
in probability, almost surely, or in $L^1$. Let $(\rho_t^M,h_t^M)$ be the
finite-width particle/boundary system and let $(\rho_t,h_t)$ be the characteristic
solution of \cref{eq:pde,eq:ode} started from $(\rho_0,h_0)$. Then the same mode
of convergence holds uniformly on $[0,T]$:
\[
  \sup_{t\in[0,T]}
  \bigl(
    \W(\rho_t^M,\rho_t)+\|h_t^M-h_t\|
  \bigr)
  \to 0.
\]
Consequently, the infinite-width predictor vector is a measurable function of the
initial environment $(\rho_0,h_0)$.
\end{corollary}

\begin{thebibliography}{9}

\bibitem{wen2020}
Yeming Wen, Dustin Tran, and Jimmy Ba.
\newblock BatchEnsemble: An alternative approach to efficient ensemble and
lifelong learning.
\newblock In \emph{International Conference on Learning Representations}, 2020.

\bibitem{jacot2018}
Arthur Jacot, Franck Gabriel, and Cl\'ement Hongler.
\newblock Neural tangent kernel: Convergence and generalization in neural
networks.
\newblock In \emph{Advances in Neural Information Processing Systems}, 2018.

\bibitem{lee2019}
Jaehoon Lee, Lechao Xiao, Samuel S. Schoenholz, Yasaman Bahri, Jascha
Sohl-Dickstein, and Jeffrey Pennington.
\newblock Wide neural networks of any depth evolve as linear models under
gradient descent.
\newblock In \emph{Advances in Neural Information Processing Systems}, 2019.

\bibitem{chizat2019}
Lena\"ic Chizat, Edouard Oyallon, and Francis Bach.
\newblock On lazy training in differentiable programming.
\newblock In \emph{Advances in Neural Information Processing Systems}, 2019.

\bibitem{mei2019}
Song Mei, Andrea Montanari, and Phan-Minh Nguyen.
\newblock A mean field view of the landscape of two-layer neural networks.
\newblock \emph{Proceedings of the National Academy of Sciences}, 115(33), 2018.

\bibitem{velikanov2022}
Maksim Velikanov, Roman V. Kail, Ivan Anokhin, Roman Vashurin, Maxim Panov,
Alexey Zaytsev, and Dmitry Yarotsky.
\newblock Embedded Ensembles: Infinite width limit and operating regimes.
\newblock In \emph{Proceedings of The 25th International Conference on
Artificial Intelligence and Statistics}, 2022.

\bibitem{gorishniy2025}
Yury Gorishniy, Akim Kotelnikov, and Artem Babenko.
\newblock TabM: Advancing tabular deep learning with parameter-efficient
ensembling.
\newblock In \emph{International Conference on Learning Representations}, 2025.

\end{thebibliography}

\appendix

\section{Abstract Mean-Field Setup and Assumptions}
\label{sec:appendix-setup}

\subsection{Abstract mean-field formulation}

At the abstract level we compress all boundary fast weights into a single symbol
$h$. Concretely, for the BatchEnsemble family above this symbol will stand for
the pair $\bar h^M=(s^{(1)},r^{(2)})$. Let $\Theta=\R^D$ denote the
single-neuron state space and let
\[
  \mathcal H=\R^{K d_h}
\]
denote the space of boundary states. Write
\[
  h=(h_1,\dots,h_K)\in\mathcal H,
  \qquad
  h_\alpha\in\R^{d_h}.
\]
Let $(\mathcal X\times\mathcal Y,\mathsf P)$ be the data space.
Fix an output dimension $d_y\in\{1,2,\dots\}$.

For each head $\alpha\in\{1,\dots,K\}$, let
\[
  \Psi_\alpha:\mathcal X\times\Theta\times\mathcal H\to\R^{d_y}
\]
be a feature map. Given a neuron law $\rho\in\PP_1(\Theta)$ and a boundary state
$h\in\mathcal H$, define
\begin{equation}
  f_\alpha(x;\rho,h)
  :=
  \int_\Theta \Psi_\alpha(x;\theta,h)\,\rho(\dd\theta).
  \label{eq:output}
\end{equation}
The head-averaged population risk is
\begin{equation}
  \mathcal R(\rho,h)
  :=
  \E_{(x,y)\sim\mathsf P}
  \Bigl[
    \frac1K\sum_{\alpha=1}^K \ell(f_\alpha(x;\rho,h),y)
  \Bigr].
  \label{eq:risk}
\end{equation}

\begin{definition}[Mean-field drifts]
Assume $\ell$ is differentiable in its first argument. Define
\begin{equation}
  B(\theta;\rho,h)
  :=
  -\E_{(x,y)\sim\mathsf P}
  \Bigl[
    \frac1K\sum_{\alpha=1}^K
    \nabla_\theta \Psi_\alpha(x;\theta,h)^\top
    \nabla_1\ell(f_\alpha(x;\rho,h),y)
  \Bigr]
  \label{eq:drift-B}
\end{equation}
and
\begin{equation}
  G_\alpha(\rho,h)
  :=
  -\E_{(x,y)\sim\mathsf P}
  \Bigl[
    \frac1K
    \int_\Theta
    \nabla_{h_\alpha}\Psi_\alpha(x;\theta,h)^\top
    \nabla_1\ell(f_\alpha(x;\rho,h),y)\,\rho(\dd\theta)
  \Bigr].
  \label{eq:drift-G}
\end{equation}
Write
\[
  G(\rho,h):=(G_\alpha(\rho,h))_{\alpha=1}^K\in\mathcal H.
\]
The fixed-$K$ mean-field system is
\begin{align}
  \partial_t \rho_t
  +
  \nabla_\theta\cdot\bigl(\rho_t B(\theta;\rho_t,h_t)\bigr)
  &=0,
  \label{eq:pde}\\
  \dot h_t &= G(\rho_t,h_t).
  \label{eq:ode}
\end{align}
\end{definition}

\begin{definition}[Characteristic solutions]
Fix $T>0$ and an initial condition $(\rho_0,h_0)\in\PP_1(\Theta)\times\mathcal H$.
A pair
\[
  (\rho_t,h_t)_{t\in[0,T]}
  \in
  C\bigl([0,T],\PP_1(\Theta)\times\mathcal H\bigr)
\]
is a characteristic solution of \cref{eq:pde,eq:ode} if there exists a flow
\[
  \Phi:[0,T]\times\Theta\to\Theta
\]
such that, for $\rho_0$-almost every $\theta$,
\[
  \dot\Phi_t(\theta)=B(\Phi_t(\theta);\rho_t,h_t),
  \qquad
  \Phi_0(\theta)=\theta,
  \qquad
  \rho_t=(\Phi_t)_\#\rho_0,
\]
while $h_t$ solves \cref{eq:ode}.
\end{definition}

\subsection{Structural assumptions}

Let $S_K$ be the permutation group on $\{1,\dots,K\}$. For $\pi\in S_K$, let
\[
  \Pi_\pi^\mathcal H(h_1,\dots,h_K)
  =
  (h_{\pi^{-1}(1)},\dots,h_{\pi^{-1}(K)}).
\]
Assume also that for each $\pi$ there is a linear isometry
$\Pi_\pi^\Theta:\Theta\to\Theta$ representing the corresponding permutation of
head-labelled coordinates inside the neuron state.

\begin{assumption}[Bounded-Lipschitz structure and permutation equivariance]
\label{ass:main}
There exist finite constants
\[
  B_\Psi,\ B_{\nabla\Psi},\ L_\Psi,\ L_{\nabla\Psi},\ L_\ell,\ C_\ell
\]
such that the following hold uniformly over all $\alpha\in\{1,\dots,K\}$ and
$x\in\mathcal X$.
\begin{enumerate}[label=(\roman*),leftmargin=2em]
  \item The maps $\Psi_\alpha(x;\cdot,\cdot)$, $\nabla_\theta\Psi_\alpha(x;\cdot,\cdot)$,
  and $\nabla_{h_\alpha}\Psi_\alpha(x;\cdot,\cdot)$ are globally bounded in
  Euclidean/operator norm by $B_\Psi$ and $B_{\nabla\Psi}$. They are also
  globally Lipschitz in $(\theta,h)$ with constants $L_\Psi$ and
  $L_{\nabla\Psi}$.
  \item $\nabla_1\ell(\cdot,y)$ is globally Lipschitz with constant $L_\ell$
  and uniformly bounded by $C_\ell$ for all $y\in\mathcal Y$.
  \item The feature maps are permutation-equivariant:
  \begin{equation}
    \Psi_\alpha(x;\Pi_\pi^\Theta\theta,\Pi_\pi^\mathcal H h)
    =
    \Psi_{\pi^{-1}(\alpha)}(x;\theta,h)
    \qquad
    \forall \pi\in S_K.
    \label{eq:equiv}
  \end{equation}
\end{enumerate}
\end{assumption}

\begin{remark}[Concrete BatchEnsemble family]
\label{rem:concrete}
For the full model \eqref{eq:model-full-matrix}--\eqref{eq:model-full}, take
\[
  \Theta_{\mathrm{BE}}
  =
  \R^{d_y}\times\R^{d_x}\times\R^K\times\R^K,
  \qquad
  \theta=(a,w,s_1^{(2)},\dots,s_K^{(2)},r_1^{(1)},\dots,r_K^{(1)}),
  \qquad
  \mathcal H_{\mathrm{BE}}=(\R^{d_x})^K\times(\R^{d_y})^K.
\]
Write a generic boundary state as
\[
  h=(u,v)\in\mathcal H_{\mathrm{BE}},
  \qquad
  u=(u_1,\dots,u_K)\in(\R^{d_x})^K,
  \qquad
  v=(v_1,\dots,v_K)\in(\R^{d_y})^K.
\]
Then \eqref{eq:model-full} is of the abstract form \eqref{eq:output} with
\begin{equation}
  \Psi_\alpha^{\mathrm{BE}}(x;\theta,h)
  :=
  (a\odot v_\alpha)\,s_\alpha^{(2)}\,
  \sigma\!\bigl(r_\alpha^{(1)}\,w^\top(u_\alpha\odot x)\bigr).
  \label{eq:psi-concrete}
\end{equation}
Fixing $(u_\alpha,v_\alpha)\equiv(\mathbf 1,\mathbf 1)$ for all heads removes
the boundary-state dependence, equivalently making the feature maps independent of
the dummy variable $h$. After composing the raw coordinates with smooth
clipping maps whose images and first derivatives are globally bounded and
permutation-equivariant, the family \eqref{eq:psi-concrete} fits
\cref{ass:main}.
\par
Accordingly, the abstract results in the main text should be read as statements
about the presence or absence of \emph{boundary fast weights}. In the full
model \eqref{eq:model-full-matrix}--\eqref{eq:model-full}, the boundary state is
the pair $(s^{(1)},r^{(2)})$; in the regime without boundary fast weights there
is no boundary state.
\end{remark}

\section{Proofs}
\label{sec:appendix-proofs}

\begin{proposition}[Output and drift Lipschitz bounds]
\label{prop:lipschitz}
Assume \cref{ass:main}. Then for every $x\in\mathcal X$ and every $\alpha$,
\begin{equation}
  \|f_\alpha(x;\rho,h)-f_\alpha(x;\mu,\bar h)\|
  \le
  L_\Psi\bigl(\W(\rho,\mu)+\|h-\bar h\|\bigr).
  \label{eq:out-lip}
\end{equation}
Moreover, there exist finite constants $L_B,L_G$ depending only on the constants
in \cref{ass:main} such that
\begin{align}
  \|B(\theta;\rho,h)-B(\eta;\mu,\bar h)\|
  &\le
  L_B\bigl(
    \|\theta-\eta\|+\W(\rho,\mu)+\|h-\bar h\|
  \bigr),
  \label{eq:B-lip}\\
  \|G(\rho,h)-G(\mu,\bar h)\|
  &\le
  L_G\bigl(\W(\rho,\mu)+\|h-\bar h\|\bigr).
  \label{eq:G-lip}
\end{align}
\end{proposition}

\begin{proof}
The output bound \eqref{eq:out-lip} follows from the dual formulation of $\W$ and
the $L_\Psi$-Lipschitz property of $\Psi_\alpha(x;\cdot,\cdot)$. For the drift
in $\theta$, decompose
\[
  B(\theta;\rho,h)-B(\eta;\mu,\bar h)=I_1+I_2,
\]
where
\begin{align*}
  I_1
  &:=
  -\E\Bigl[
    \frac1K\sum_{\alpha=1}^K
    \nabla_\theta\Psi_\alpha(x;\theta,h)^\top
    \bigl(
      \nabla_1\ell(f_\alpha(x;\rho,h),y)
      -
      \nabla_1\ell(f_\alpha(x;\mu,\bar h),y)
    \bigr)
  \Bigr],\\
  I_2
  &:=
  -\E\Bigl[
    \frac1K\sum_{\alpha=1}^K
    \bigl(
      \nabla_\theta\Psi_\alpha(x;\theta,h)
      -
      \nabla_\theta\Psi_\alpha(x;\eta,\bar h)
    \bigr)^\top
    \nabla_1\ell(f_\alpha(x;\mu,\bar h),y)
  \Bigr].
\end{align*}
Using \eqref{eq:out-lip}, the Lipschitz property of $\nabla_1\ell$, and the
bounds in \cref{ass:main} yields \eqref{eq:B-lip}. The proof of \eqref{eq:G-lip}
is identical, replacing $\nabla_\theta\Psi_\alpha$ with
$\nabla_{h_\alpha}\Psi_\alpha$.
\end{proof}

\begin{proof}[Proof of \cref{thm:deterministic}]
By \cref{prop:lipschitz}, the coupled vector field in
\cref{eq:pde,eq:ode} is globally Lipschitz in $(\theta,\rho,h)$ on
finite-horizon balls, so standard Picard iteration yields existence and
uniqueness of a characteristic solution.

For stability, let $\Phi_t$ and $\bar\Phi_t$ be the characteristic flows of the
two solutions. Start from an optimal coupling of $\rho_0$ and $\mu_0$, transport
it by $(\Phi_t,\bar\Phi_t)$, and combine the resulting estimate with
\eqref{eq:B-lip} and \eqref{eq:G-lip}. Gronwall then yields
\eqref{eq:det-stability}.

For particle convergence, note that $(\rho_t^M,h_t^M)$ is itself a characteristic
solution of \cref{eq:pde,eq:ode} with initial data $(\rho_0^M,h_0^M)$. Applying
\eqref{eq:det-stability} with $(\mu_t,\bar h_t)=(\rho_t^M,h_t^M)$ gives
\eqref{eq:particle-conv}, and the predictor convergence follows from
\eqref{eq:out-lip}.
\end{proof}

\begin{proof}[Proof of \cref{cor:random}]
This is an immediate consequence of \eqref{eq:particle-conv}.
\end{proof}

\begin{proof}[Proof of \cref{cor:boundaryfree}]
If $\Psi_\alpha$ does not depend on $h$, then $\nabla_{h_\alpha}\Psi_\alpha=0$
and hence $G\equiv 0$. The ODE component decouples and can be discarded.
\end{proof}

\begin{proof}[Proof of \cref{cor:boundaryfree-det}]
This is \cref{cor:random} with no boundary fast weights present.
\end{proof}

\begin{proof}[Proof of \cref{prop:collapse}]
Define $\tilde\rho_t:=(\Pi_\pi^\Theta)_\#\rho_t$. Because the feature maps do not
depend on $h$, \eqref{eq:equiv} reduces to
\[
  \Psi_\alpha^{\mathrm{free}}(x;\Pi_\pi^\Theta\theta)
  =
  \Psi_{\pi^{-1}(\alpha)}^{\mathrm{free}}(x;\theta).
\]
Hence
\begin{align*}
  f_\alpha(x;\tilde\rho_t)
  &=
  \int \Psi_\alpha^{\mathrm{free}}(x;\theta)\,\tilde\rho_t(\dd\theta)\\
  &=
  \int \Psi_\alpha^{\mathrm{free}}(x;\Pi_\pi^\Theta\eta)\,\rho_t(\dd\eta)\\
  &=
  \int \Psi_{\pi^{-1}(\alpha)}^{\mathrm{free}}(x;\eta)\,\rho_t(\dd\eta)\\
  &=
  f_{\pi^{-1}(\alpha)}(x;\rho_t).
\end{align*}
Applying the same identity to the drift shows
\[
  B(\Pi_\pi^\Theta\theta;\tilde\rho_t)
  =
  \Pi_\pi^\Theta B(\theta;\rho_t).
\]
Therefore $\tilde\rho_t$ solves the same PDE as $\rho_t$ with the same initial
law. Uniqueness gives $\tilde\rho_t=\rho_t$. The collapse identity
\eqref{eq:collapse} then follows by choosing a permutation sending any head
index to any other one.
\end{proof}

\begin{proof}[Proof of \cref{cor:randomness}]
This is immediate from \cref{cor:random}, because the limit trajectory is a
measurable function of $(\rho_0,h_0)$.
\end{proof}

\end{document}
